% Shared layout for the two Lab0 translations: Chinese two-column conference style.
\usepackage[a4paper,left=1.9cm,right=1.9cm,top=2.4cm,bottom=2.2cm,
  columnsep=0.7cm,headsep=0.5cm]{geometry}
\usepackage{amsmath,amssymb,bm}
\setmainfont{Times New Roman}
\setsansfont{Helvetica}
\setmonofont{Menlo}[Scale=0.85]
% macOS system CJK fonts; Kai/Fang come from TeX Live's Fandol set.
\setCJKmainfont{Songti SC}[BoldFont={Heiti SC},ItalicFont={FandolKai-Regular.otf}]
\setCJKsansfont{Heiti SC}
\setCJKmonofont{Heiti SC}
\setCJKfamilyfont{zhhei}{Heiti SC}
\setCJKfamilyfont{zhkai}{FandolKai-Regular.otf}
\setCJKfamilyfont{zhfs}{FandolFang-Regular.otf}
\providecommand{\heiti}{\CJKfamily{zhhei}}
\providecommand{\kaishu}{\CJKfamily{zhkai}}
\providecommand{\fangsong}{\CJKfamily{zhfs}}
\usepackage{newtxmath}
% fontspec's text fonts have no OT1 shapes; point math alphabets at newtx's Times.
\DeclareMathAlphabet{\mathrm}{OT1}{ntxtlf}{m}{n}
\DeclareMathAlphabet{\mathbf}{OT1}{ntxtlf}{b}{n}
\DeclareMathAlphabet{\mathtt}{OT1}{txtt}{m}{n}
\DeclareMathAlphabet{\mathsf}{OT1}{phv}{m}{n}
\SetSymbolFont{operators}{normal}{OT1}{ntxtlf}{m}{n}
\SetSymbolFont{operators}{bold}{OT1}{ntxtlf}{b}{n}
\usepackage{graphicx,booktabs,multirow,array}
\usepackage[font=small,labelsep=quad,skip=4pt]{caption}
\usepackage{enumitem}
\usepackage{fancyvrb}
\usepackage{fancyhdr}
\usepackage{xurl}
\usepackage[hidelinks]{hyperref}
\graphicspath{{figures/}}
\urlstyle{same}

\ctexset{
  section={format=\heiti\zihao{-4}\raggedright,beforeskip=1.2ex plus .3ex,
    afterskip=.8ex plus .2ex,aftername=\quad},
  subsection={format=\heiti\zihao{5}\raggedright,beforeskip=.9ex plus .2ex,
    afterskip=.5ex plus .1ex,aftername=\quad},
}
\setlist{nosep,leftmargin=2em}
\setlength{\parskip}{0pt}
\linespread{1.18}
\setlength{\abovedisplayskip}{5pt plus 2pt minus 2pt}
\setlength{\belowdisplayskip}{5pt plus 2pt minus 2pt}
\setlength{\textfloatsep}{10pt plus 2pt minus 3pt}
\setlength{\dbltextfloatsep}{10pt plus 2pt minus 3pt}
\emergencystretch=1em

% Bold run-in heading used by both papers ("Photogrammetric calibration." etc.).
\newcommand{\runin}[1]{\textbf{\heiti #1}\quad}
% Translator's note in a footnote.
\newcommand{\yizhe}[1]{\footnote{\kaishu 译者注：#1}}

\pagestyle{fancy}
\fancyhf{}
\renewcommand{\headrulewidth}{0.4pt}
\fancyfoot[C]{\small\thepage}
